\section{Landasan Matematis}
\label{sec:landasan-elgamal}

Keamanan algoritma ElGamal terletak pada sulitnya menghitung logaritma diskret.
Sebelum membahas algoritmanya, bagian ini menjelaskan akar primitif dan
persoalan logaritma diskret. Keduanya dipakai pada seluruh prosedur ElGamal.

\subsection{Akar Primitif}

Misalkan $n$ adalah bilangan bulat. Bilangan $a$ disebut \textbf{akar primitif}
dari $n$ bila perpangkatan
\[ a, a^{2}, \dots, a^{\varphi(n)} \pmod{n} \]
menghasilkan nilai yang berbeda dan semuanya relatif prima dengan $n$.

Secara khusus, misalkan $p$ adalah bilangan prima. Bilangan $a$ disebut akar
primitif dari $p$ bila perpangkatan
\[ a, a^{2}, \dots, a^{p-1} \pmod{p} \]
menghasilkan nilai-nilai yang berbeda. Hal ini terjadi karena untuk $p$ prima
berlaku $\varphi(p) = p - 1$.

\subsection{Contoh pada Modulus Kecil}

Misalkan $p = 7$. Bilangan $a = 3$ adalah akar primitif dari $7$ karena
\[
\begin{aligned}
3^{1} \pmod 7 &= 3; & 3^{2} \pmod 7 &= 2; & 3^{3} \pmod 7 &= 6 \\
3^{4} \pmod 7 &= 4; & 3^{5} \pmod 7 &= 5; & 3^{6} \pmod 7 &= 1 .
\end{aligned}
\]
Jadi, semua perpangkatan dari $3$ menghasilkan nilai yang berbeda, yaitu
$(3, 2, 6, 4, 5, 1)$. Semua bilangan di dalam modulus $7$ terjadi satu kali.

Perpangkatan berikutnya, $3^{7} \pmod 7$, $3^{8} \pmod 7$, dan seterusnya, akan
kembali berulang menghasilkan nilai-nilai tersebut. Panjang satu siklus tidak
lebih dari $7 - 1 = 6$. Secara umum, untuk $p$ bilangan prima, panjang siklus
tidak lebih dari $p - 1$.

\begin{table}[htbp]
    \centering
    \begin{tabular}{@{}l c c c c c c c@{}}
    \toprule
    $i$ & 1 & 2 & 3 & 4 & 5 & 6 & 7 \\
    \midrule
    $3^{i} \bmod 7$ & 3 & 2 & 6 & 4 & 5 & \textbf{1} & 3 \\
    $2^{i} \bmod 7$ & 2 & 4 & \textbf{1} & 2 & 4 & \textbf{1} & 2 \\
    \bottomrule
    \end{tabular}
    \caption{Siklus perpangkatan modulo 7. Baris $a = 3$ memuat seluruh residu
    bukan-nol tepat sekali, sedangkan baris $a = 2$ hanya memuat tiga nilai
    berbeda dan berulang dengan periode 3}
    \label{tab:akar-primitif-7}
\end{table}

Sebaliknya, $a = 2$ bukan akar primitif dari $7$ karena
\[
\begin{aligned}
2^{1} \pmod{7} &= 2; & 2^{2} \pmod{7} &= 4; & 2^{3} \pmod{7} &= 1 \\
2^{4} \pmod{7} &= 2; & 2^{5} \pmod{7} &= 4; & 2^{6} \pmod{7} &= 1 .
\end{aligned}
\]
Nilai-nilai yang dihasilkan dari $2^{1}, 2^{2}, \dots, 2^{6}$ tidak semuanya
berbeda dan tidak mencakup semua nilai di dalam modulus $7$. Untuk menemukan
semua akar primitif dari $p$, kita harus mencoba semua bilangan bulat dari
$2, 3, \dots$

\subsection{Persoalan Logaritma Diskret}

Jika $a$ adalah akar primitif dari bilangan prima $p$, maka untuk bilangan bulat
$b$ kita dapat menemukan pangkat $x$ sedemikian sehingga
\begin{equation}
b \equiv a^{x} \pmod{p}, \qquad 0 \le x \le p-1 .
\label{eq:dlp-elgamal}
\end{equation}
Pangkat $x$ disebut \textbf{logaritma diskret} dari $b$ untuk basis $a$ modulo
$p$.

Dalam \textbf{persoalan logaritma diskret}, diberikan $b \equiv a^{x} \pmod{p}$,
carilah $x$ yang memenuhi kekongruenan tersebut.

Sebagai contoh, $7$ adalah akar primitif dari bilangan prima $p = 41$. Kita
mencari $x$ sedemikian sehingga
\[ 15 \equiv 7^{x} \pmod{41} . \]
Jawabannya adalah $x = 3$ karena
\[ 7^{3} = 343 \equiv 15 \pmod{41} . \]
Persoalan inilah yang menjadi landasan keamanan algoritma ElGamal.
